\chapter{1989年广东卷（理）}

\section{单选题}

\begin{problem}
直线 \(x+\sqrt{3}y-3=0\) 的倾斜角是（\(\quad\)）

\choices
    {\(\frac{\pi}{6}\)}
    {\(\frac{\pi}{3}\)}
    {\(\frac{5\pi}{6}\)}
    {\(\frac{2\pi}{3}\)}
\end{problem}

\begin{answer}
C
\end{answer}

\begin{solution}
将直线方程化为斜截式，得
\(y=-\frac{1}{\sqrt{3}}x+\sqrt{3}\)，所以直线的斜率为 \(\tan\alpha=-\frac{1}{\sqrt{3}}\). 直线的倾斜角满足 \(0\leqslant\alpha<\pi\)，故 \(\alpha=\frac{5\pi}{6}\)，应选 C.
\end{solution}

\begin{problem}
二面角是指（\(\quad\)）

\choices
    {两个平面所组成的角}
    {经过同一直线的两个平面所组成的图形}
    {从一条直线出发的两个半平面所组成的图形}
    {两个平面所夹的不大于 \(90^\circ\) 的角}
\end{problem}

\begin{answer}
C
\end{answer}

\begin{solution}
二面角由同一条棱出发的两个半平面组成，两个半平面的公共边就是二面角的棱. 因此题中准确的定义是“从一条直线出发的两个半平面所组成的图形”，应选 C.
\end{solution}

\begin{problem}
函数 \(f(x)=x^{\frac{3}{5}}\)（\(\quad\)）

\choices
    {是奇函数而不是偶函数}
    {是偶函数而不是奇函数}
    {既是奇函数又是偶函数}
    {既不是奇函数又不是偶函数}
\end{problem}

\begin{answer}
A
\end{answer}

\begin{solution}
函数的定义域为 \(\mathbb{R}\)，且对任意 \(x\in\mathbb{R}\)，有 \(f(-x)=(-x)^{\frac{3}{5}}=-x^{\frac{3}{5}}=-f(x)\)，所以它是奇函数. 又例如 \(f(1)=1\)，\(f(-1)=-1\)，二者不相等，故不是偶函数，应选 A.
\end{solution}

\begin{problem}
如果 \(\sin\left( \pi+A \right)=-\frac{1}{2}\)，那么 \(\cos\left( \frac{3\pi}{2}-A \right)=\)（\(\quad\)）

\choices
    {\(-\frac{1}{2}\)}
    {\(\frac{1}{2}\)}
    {\(-\frac{\sqrt{3}}{2}\)}
    {\(\frac{\sqrt{3}}{2}\)}
\end{problem}

\begin{answer}
A
\end{answer}

\begin{solution}
由诱导公式，\(\sin(\pi+A)=-\sin A\)，故 \(-\sin A=-\frac12\)，从而 \(\sin A=\frac12\). 又
\(\cos\left(\frac{3\pi}{2}-A\right)=\cos\frac{3\pi}{2}\cos A+\sin\frac{3\pi}{2}\sin A=-\sin A=-\frac12\)，应选 A.
\end{solution}

\begin{problem}
函数 \(y=4^{-x}\) 的反函数是（\(\quad\)）

\choices
    {\(y=\log_4\left( -x \right)\)}
    {\(y=-\log_4 x\)}
    {\(y=\log_4 x\)}
    {\(y=4^x\)}
\end{problem}

\begin{answer}
B
\end{answer}

\begin{solution}
设原函数中自变量为 \(t\)，则 \(y=4^{-t}\) 等价于 \(\log_4 y=-t\)，即 \(t=-\log_4 y\). 交换自变量与函数值的字母，反函数为 \(y=-\log_4 x\)，其定义域为 \(x>0\)，应选 B.
\end{solution}

\begin{problem}
设 \(A\) 是直角坐标平面上所有的点所组成的集合，如果由 \(A\) 到 \(A\) 的一一对应映射 \(f\)，使象集合的元素 \(\left( y-1,x+2 \right)\) 和原象集合的元素 \(\left( x,y \right)\) 对应，那么，象点 \((3,-4)\) 的原象是点（\(\quad\)）

\choices
    {\((-5,5)\)}
    {\((4,-6)\)}
    {\((2,-2)\)}
    {\((-6,4)\)}
\end{problem}

\begin{answer}
D
\end{answer}

\begin{solution}
映射把原象点 \((x,y)\) 变为象点 \((y-1,x+2)\). 若象点为 \((3,-4)\)，则
\(y-1=3\)，\(x+2=-4\)，所以 \(y=4\)，\(x=-6\). 因而原象点为 \((-6,4)\)，应选 D.
\end{solution}

\begin{problem}
在等比数列 \(\{a_n\}\) 中，如果 \(a_6=6\)，\(a_9=9\)，那么 \(a_3=\)（\(\quad\)）

\choices
    {\(4\)}
    {\(\frac{3}{2}\)}
    {\(\frac{16}{9}\)}
    {\(3\)}
\end{problem}

\begin{answer}
A
\end{answer}

\begin{solution}
设等比数列的公比为 \(q\). 由 \(a_9=a_6q^3\)，得 \(9=6q^3\)，即 \(q^3=\frac32\). 又 \(a_6=a_3q^3\)，所以
\(a_3=\frac{a_6}{q^3}=\frac{6}{3/2}=4\)，应选 A.
\end{solution}

\begin{problem}
如果抛物线的顶点在原点，对称轴对 \(x\) 轴，焦点在直线 \(3x-4y-12=0\) 上，那么抛物线的方程是（\(\quad\)）

\choices
    {\(y^2=-16x\)}
    {\(y^2=12x\)}
    {\(y^2=16x\)}
    {\(y^2=-12x\)}
\end{problem}

\begin{answer}
C
\end{answer}

\begin{solution}
顶点在原点且对称轴为 \(x\) 轴的抛物线可写为 \(y^2=4ax\)，焦点为 \((a,0)\). 焦点在直线 \(3x-4y-12=0\) 上，代入 \((a,0)\) 得 \(3a-12=0\)，所以 \(a=4\). 因而抛物线方程为 \(y^2=16x\)，应选 C.
\end{solution}

\begin{problem}
如果椭圆的两条准线间的距离是这个椭圆的焦距的两倍，那么这个椭圆的离心率为（\(\quad\)）

\choices
    {\(\frac{1}{4}\)}
    {\(\frac{1}{2}\)}
    {\(\frac{\sqrt{2}}{2}\)}
    {\(\frac{\sqrt{2}}{4}\)}
\end{problem}

\begin{answer}
C
\end{answer}

\begin{solution}
设椭圆的半长轴、半焦距分别为 \(a\)、\(c\)，离心率为 \(e=\frac ca\). 两条准线间的距离为 \(\frac{2a}{e}\)，焦距为 \(2c=2ae\). 由题意
\(\frac{2a}{e}=2\cdot2ae\)，故 \(e^2=\frac12\). 因为 \(0<e<1\)，所以 \(e=\frac{\sqrt2}{2}\)，应选 C.
\end{solution}

\begin{problem}
如果双曲线 \(\frac{x^2}{9}-y^2=1\) 的两个焦点为 \(F_1\)，\(F_2\)，\(A\) 是该双曲线上一点，且 \(\left| AF_1 \right|=5\)，那么 \(\left| AF_2 \right|=\)（\(\quad\)）

\choices
    {\(5+\sqrt{10}\)}
    {\(5+2\sqrt{10}\)}
    {\(8\)}
    {\(11\)}
\end{problem}

\begin{answer}
D
\end{answer}

\begin{solution}
双曲线的半实轴为 \(a=3\)，半虚轴为 \(b=1\)，故半焦距 \(c=\sqrt{a^2+b^2}=\sqrt{10}\). 双曲线上一点到两焦点的距离之差的绝对值恒为 \(2a=6\). 已知 \(\lvert AF_1\rvert=5\)，另一段距离只能满足 \(\lvert AF_2\rvert=5+6=11\)，应选 D.
\end{solution}

\begin{problem}
如果 \(l\) 和 \(n\) 是异面直线，那么和 \(l\)，\(n\) 都垂直的直线（\(\quad\)）

\choices
    {不一定存在}
    {总共只有一条}
    {总共可能有一条，也可能有两条}
    {必然有无穷多条}
\end{problem}

\begin{answer}
D
\end{answer}

\begin{solution}
异面直线 \(l\)、\(n\) 的方向向量不平行，因此存在一个同时垂直于二者方向的方向. 取该方向作一条直线，便得到同时垂直于 \(l\)、\(n\) 的直线；将这条直线平移到空间中任意位置，仍与两直线垂直，所以这样的直线有无穷多条. 应选 D.
\end{solution}

\begin{problem}
“两底面直径之差等于母线长”的圆台（\(\quad\)）

\choices
    {是不存在的}
    {其母线与下底面必成 \(60^\circ\) 角}
    {其母线与下底面所成的角不是定值}
    {其高与母线必成 \(60^\circ\) 角}
\end{problem}

\begin{answer}
B
\end{answer}

\begin{solution}
取过圆台轴线的截面，得到等腰梯形. 设圆台的母线长为 \(l\)，上下底面半径分别为 \(R\)、\(r\)，母线与下底面所成的角为 \(\alpha\). 题设“两底面直径之差等于母线长”给出 \(2R-2r=l\)，即母线在底面内的投影长度 \(R-r=\frac l2\). 因而
\(\cos\alpha=\frac{R-r}{l}=\frac12\)，且 \(0<\alpha<\frac\pi2\)，所以 \(\alpha=60^\circ\)，应选 B.
\end{solution}

\begin{problem}
侧面都是直角三角形的正三棱锥，底面边长为 \(a\) 时，该三棱锥的全面积是（\(\quad\)）

\choices
    {\(\frac{3+\sqrt{3}}{4}a^2\)}
    {\(\frac{3}{4}a^2\)}
    {\(\frac{3+\sqrt{3}}{2}a^2\)}
    {\(\frac{6+\sqrt{3}}{4}a^2\)}
\end{problem}

\begin{answer}
A
\end{answer}

\begin{solution}
设正三棱锥的顶点为 \(S\)，底面为正三角形 \(ABC\). 每个侧面是等腰三角形；若其为直角三角形，直角不能在底面顶点，只能在顶点 \(S\). 因而每个侧面为等腰直角三角形，底边 \(a\) 是斜边，每个侧面的面积为 \(\frac12\left(\frac{a}{\sqrt2}\right)^2=\frac{a^2}{4}\). 三个侧面积之和为 \(\frac{3a^2}{4}\)，底面面积为 \(\frac{\sqrt3}{4}a^2\)，所以全面积为 \(\frac{3+\sqrt3}{4}a^2\)，应选 A.
\end{solution}

\begin{problem}
设 \(\{a_n\}\) 是公差为 \(-2\) 的等差数列，如果 \(a_1+a_4+a_7+\cdots+a_{97}=50\)，那么 \(a_3+a_6+a_9+\cdots+a_{99}=\)（\(\quad\)）

\choices
    {\(-182\)}
    {\(-78\)}
    {\(-148\)}
    {\(-82\)}
\end{problem}

\begin{answer}
D
\end{answer}

\begin{solution}
设 \(a_1=A\). 因为公差为 \(-2\)，所以
\(a_{97}=A-2\cdot96=A-192\). 等差数列 \(a_1\)，\(a_4\)，\(\ldots\)，\(a_{97}\) 有 \(33\) 项，故
\(33\cdot\frac{a_1+a_{97}}2=33(A-96)=50\).
目标数列 \(a_3\)，\(a_6\)，\(\ldots\)，\(a_{99}\) 也有 \(33\) 项，且 \(a_3=A-4\)，\(a_{99}=A-196\)，所以其和为
\(33\cdot\frac{(A-4)+(A-196)}2=33(A-100)=50-132=-82\)，应选 D.
\end{solution}

\begin{problem}
如果直线 \(l_1\)，\(l_2\) 的斜率分别是二次方程 \(x^2-4x+1=0\) 的两根，那么 \(l_1\)，\(l_2\) 所成的角是（\(\quad\)）

\choices
    {\(\frac{\pi}{3}\)}
    {\(\frac{\pi}{4}\)}
    {\(\frac{\pi}{6}\)}
    {\(\frac{\pi}{8}\)}
\end{problem}

\begin{answer}
A
\end{answer}

\begin{solution}
方程 \(x^2-4x+1=0\) 的两根为 \(m_1=2+\sqrt3\)、\(m_2=2-\sqrt3\). 于是 \(m_1m_2=1\)，且两直线的锐角 \(\varphi\) 满足
\(\tan\varphi=\left|\frac{m_1-m_2}{1+m_1m_2}\right|=\frac{2\sqrt3}{2}=\sqrt3\). 因而 \(\varphi=\frac\pi3\)，应选 A.
\end{solution}

\begin{problem}
在极坐标系中，如果等边三角形 \(ABC\) 的两个顶点是 \(A\left( 2,\frac{\pi}{4} \right)\)，\(B\left( 2,\frac{5\pi}{4} \right)\)，那么顶点 \(C\) 的坐标可能是（\(\quad\)）

\choices
    {\(\left( 4,\frac{3\pi}{4} \right)\)}
    {\(\left( 2\sqrt{3},\frac{3\pi}{4} \right)\)}
    {\(\left( 2\sqrt{3},\pi \right)\)}
    {\((3,\pi)\)}
\end{problem}

\begin{answer}
B
\end{answer}

\begin{solution}
极坐标点 \(A\)、\(B\) 的直角坐标分别为 \((\sqrt2,\sqrt2)\)、\((-\sqrt2,-\sqrt2)\)，所以 \(AB=4\)，且其中点为原点. 等边三角形的高为 \(2\sqrt3\)，方向垂直于 \(AB\)，因此一个可能的顶点为极坐标 \(\left(2\sqrt3,\frac{3\pi}{4}\right)\). 代入直角坐标可得该点为 \((-\sqrt6,\sqrt6)\)，到 \(A\)、\(B\) 的距离均为 \(4\)，应选 B.
\end{solution}

\begin{problem}
不等式 \(\sqrt{4-x^2}+\frac{\left| x \right|}{x}\geqslant 0\) 的解集是（\(\quad\)）

\choices
    {\(\left\{ x\mid -2\leqslant x\leqslant 2 \right\}\)}
    {\(\left\{ x\mid -\sqrt{3}\leqslant x<0\text{或}0<x\leqslant 2 \right\}\)}
    {\(\left\{ x\mid -2\leqslant x<0\text{或}0<x\leqslant 2 \right\}\)}
    {\(\left\{ x\mid -\sqrt{3}\leqslant x<0\text{或}0<x\leqslant \sqrt{3} \right\}\)}
\end{problem}

\begin{answer}
B
\end{answer}

\begin{solution}
根式有意义要求 \(4-x^2\geqslant0\)，即 \(-2\leqslant x\leqslant2\)，且 \(x\neq0\). 当 \(x>0\) 时，\(\frac{|x|}{x}=1\)，不等式显然成立，得到 \(0<x\leqslant2\). 当 \(x<0\) 时，\(\frac{|x|}{x}=-1\)，于是
\(\sqrt{4-x^2}\geqslant1\)，等价于 \(4-x^2\geqslant1\)，从而 \(-\sqrt3\leqslant x<0\). 所以解集为 \(\left[-\sqrt3,0\right)\cup(0,2]\)，应选 B.
\end{solution}

\begin{problem}
如果函数 \(f(x)=x^2+2(a-1)x+2\) 在区间 \((-\infty,4]\) 上是减函数，那么实数 \(a\) 的取值范围是（\(\quad\)）

\choices
    {\(a\geqslant -3\)}
    {\(a\leqslant -3\)}
    {\(a\leqslant 5\)}
    {\(a\geqslant 3\)}
\end{problem}

\begin{answer}
B
\end{answer}

\begin{solution}
函数可写为 \(f(x)=\left(x+a-1\right)^2+2-(a-1)^2\)，其对称轴为 \(x=1-a\). 抛物线开口向上，因此在 \((-\infty,4]\) 上递减，当且仅当该区间的右端点不在对称轴右侧，即 \(4\leqslant1-a\). 故 \(a\leqslant-3\)，应选 B.
\end{solution}

\begin{problem}
如果 \(\tan \frac{A}{2}=\frac{m}{n}\)，那么 \(m\cos A-n\sin A=\)（\(\quad\)）

\choices
    {\(n\)}
    {\(-n\)}
    {\(-m\)}
    {\(m\)}
\end{problem}

\begin{answer}
C
\end{answer}

\begin{solution}
由半角正切公式，令 \(t=\tan\frac A2=\frac mn\)，则
\(\cos A=\frac{1-t^2}{1+t^2}=\frac{n^2-m^2}{n^2+m^2}\)，\(\sin A=\frac{2t}{1+t^2}=\frac{2mn}{n^2+m^2}\). 因而
\[
m\cos A-n\sin A
=\frac{m(n^2-m^2)-2mn^2}{n^2+m^2}
=-m
\]
应选 C.
\end{solution}

\begin{problem}
如果 \(0<m<b<a\)，那么（\(\quad\)）

\choices
    {\(\cos \frac{b+m}{a+m}<\cos \frac{b}{a}<\cos \frac{b-m}{a-m}\)}
    {\(\cos \frac{b}{a}<\cos \frac{b-m}{a-m}<\cos \frac{b+m}{a+m}\)}
    {\(\cos \frac{b-m}{a-m}<\cos \frac{b}{a}<\cos \frac{b+m}{a+m}\)}
    {\(\cos \frac{b+m}{a+m}<\cos \frac{b-m}{a-m}<\cos \frac{b}{a}\)}
\end{problem}

\begin{answer}
A
\end{answer}

\begin{solution}
由 \(0<m<b<a\)，有
\(\frac{b+m}{a+m}-\frac ba=\frac{m(a-b)}{a(a+m)}>0\)，\(\frac ba-\frac{b-m}{a-m}=\frac{m(a-b)}{a(a-m)}>0\).
所以 \(\frac{b-m}{a-m}<\frac ba<\frac{b+m}{a+m}\). 三个比值都在 \((0,1)\) 内，余弦函数在该区间上递减，故
\(\cos\frac{b+m}{a+m}<\cos\frac ba<\cos\frac{b-m}{a-m}\)，应选 A.
\end{solution}

\begin{problem}
如果 \(\theta\) 是第三象限的角，且满足 \(\sqrt{1+\sin \theta}=\cos \frac{\theta}{2}+\sin \frac{\theta}{2}\)，那么 \(\frac{\theta}{2}\) 是（\(\quad\)）

\choices
    {第四象限的角}
    {第三象限的角}
    {第二象限的角}
    {第一象限的角}
\end{problem}

\begin{answer}
C
\end{answer}

\begin{solution}
令 \(t=\frac\theta2\). 因为 \(1+\sin\theta=1+2\sin t\cos t=(\sin t+\cos t)^2\)，原条件等价于
\(\lvert\sin t+\cos t\rvert=\sin t+\cos t\)，即 \(\sin t+\cos t\geqslant0\). 又 \(\theta\) 为第三象限的角，所以 \(t\) 位于第二象限的相应范围内；在该范围内 \(\sin t+\cos t\geqslant0\). 因而 \(\frac\theta2\) 是第二象限的角，应选 C.
\end{solution}

\begin{problem}
如果 \(x\)，\(y\) 是实数，那么“\(x\neq y\)”是“\(\cos x\neq \cos y\)”的（\(\quad\)）

\choices
    {充分而且必要的条件}
    {充分但不必要的条件}
    {必要但不充分的条件}
    {既不充分也不必要的条件}
\end{problem}

\begin{answer}
C
\end{answer}

\begin{solution}
若 \(\cos x\neq\cos y\)，则不可能有 \(x=y\)，所以 \(x\neq y\) 是 \(\cos x\neq\cos y\) 的必要条件. 但它不是充分条件，例如取 \(x=0\)、\(y=2\pi\)，有 \(x\neq y\)，而 \(\cos x=\cos y=1\). 因此该条件必要但不充分，应选 C.
\end{solution}

\begin{problem}
方程 \(\arcsin x+\arccos y=\frac{\pi}{2}\) 所表示的曲线是（\(\quad\)）

\choices
    {一个半圆弧}
    {一个圆周}
    {一条直线}
    {一条线段}
\end{problem}

\begin{answer}
D
\end{answer}

\begin{solution}
反正弦函数的值域为 \(\left[-\frac\pi2,\frac\pi2\right]\)，反余弦函数的值域为 \([0,\pi]\)，并且在各自的主值范围内有 \(\arccos y=\frac\pi2-\arcsin y\). 因而原方程等价于
\(\arcsin x=\arcsin y\)，从而 \(x=y\)，且 \(-1\leqslant x\leqslant1\). 所表示的曲线是直线 \(y=x\) 上从 \((-1,-1)\) 到 \((1,1)\) 的线段，应选 D.
\end{solution}

\begin{problem}
如果 \(0<a<1\)，\(0<x\leqslant y<1\)，且 \(\log_a x\log_a y=1\)，那么 \(xy\)（\(\quad\)）

\choices
    {无最大值也无最小值}
    {无最大值而有最小值}
    {有最大值而无最小值}
    {有最大值也有最小值}
\end{problem}

\begin{answer}
C
\end{answer}

\begin{solution}
令 \(u=\log_a x\)、\(v=\log_a y\). 因为 \(0<a<1\) 且 \(0<x\leqslant y<1\)，所以 \(u\geqslant v>0\)，又 \(uv=1\). 于是可令 \(u\geqslant1\)，\(v=\frac1u\)，从而
\(xy=a^{u+v}=a^{u+1/u}\). 对于 \(1\leqslant u_1<u_2\)，有
\(\left(u_2+\frac1{u_2}\right)-\left(u_1+\frac1{u_1}\right)=(u_2-u_1)\left(1-\frac1{u_1u_2}\right)>0\)，所以 \(u+\frac1u\) 在 \([1,+\infty)\) 上递增，最小值为 \(2\)，且随着 \(u\to+\infty\) 趋于 \(+\infty\). 由于 \(0<a<1\)，故 \(xy\) 在 \(u=1\) 时取得最大值 \(a^2\)，但当 \(u\to+\infty\) 时趋于 \(0\)，永远不能取到 \(0\). 所以 \(xy\) 有最大值而无最小值，应选 C.
\end{solution}

\begin{problem}
设 \(k\) 是正实数，如果方程 \(kxy+x^2-x+4y-6=0\) 表示两条直线，那么，它们的图形是（\(\quad\)）

\choices
    {\choicebitmap[width=0.15\paperwidth]{img/1989/guangdong_science/q25_optA.png}}
    {\choicebitmap[width=0.15\paperwidth]{img/1989/guangdong_science/q25_optB.png}}
    {\choicebitmap[width=0.15\paperwidth]{img/1989/guangdong_science/q25_optC.png}}
    {\choicebitmap[width=0.15\paperwidth]{img/1989/guangdong_science/q25_optD.png}}
\end{problem}

\begin{answer}
B
\end{answer}

\begin{solution}
将方程按 \(y\) 分组，得
\(y(kx+4)+(x-3)(x+2)=0\). 若它表示两条直线且含有 \(x^2\) 项，则其中一条直线必须为竖直直线，设分解为
\((x+\lambda)(x+ky+\mu)=0\). 比较系数，得到
\(\lambda k=4\)，\(\lambda\mu=-6\)，\(\lambda+\mu=-1\). 由前两式得 \(\lambda=\frac4k\)、\(\mu=-\frac{3k}{2}\)，代入第三式：
\(\frac4k-\frac{3k}{2}=-1\)，所以 \(3k^2-2k-8=0\). 因为 \(k>0\)，得 \(k=2\)，方程分解为
\((x+2)(x+2y-3)=0\).

两条直线分别为 \(x=-2\) 和 \(y=-\frac12x+\frac32\)，后者与坐标轴的截距为 \(3\)、\(\frac32\)，图形对应选项 B.
\end{solution}

\section{填空题}

\begin{problem}
函数 \(y=\sqrt{x^2-49}\) 的值域是\fillinblank{}.
\end{problem}

\begin{answer}
\(\left[0,+\infty\right)\)
\end{answer}

\begin{solution}
根式有意义的条件为 \(x^2-49\geqslant0\)，即 \(\lvert x\rvert\geqslant7\). 对任意 \(t\geqslant0\)，取 \(x=\sqrt{t^2+49}\)，则 \(\sqrt{x^2-49}=t\). 因而函数值可以取到所有非负实数，值域为 \(\left[0,+\infty\right)\).
\end{solution}

\begin{problem}
\(\sin 75^\circ-\sin 15^\circ\) 的值是\fillinblank{}.
\end{problem}

\begin{answer}
\(\frac{\sqrt{2}}{2}\)
\end{answer}

\begin{solution}
利用正弦差公式，得
\[
\sin75^\circ-\sin15^\circ
=2\cos45^\circ\sin30^\circ
=2\cdot\frac{\sqrt2}{2}\cdot\frac12
=\frac{\sqrt2}{2}
\]
所以填 \(\frac{\sqrt2}{2}\).
\end{solution}

\begin{problem}
\(\displaystyle\lim_{n\to\infty}\left( 1+\frac{1}{3}+\frac{1}{3^2}+\cdots+\frac{1}{3^n} \right)\) 的值是\fillinblank{}.
\end{problem}

\begin{answer}
\(\frac{3}{2}\)
\end{answer}

\begin{solution}
括号内是首项为 \(1\)、公比为 \(\frac13\) 的等比数列前 \(n+1\) 项之和，因此
\[
1+\frac13+\cdots+\frac1{3^n}
=\frac{1-(1/3)^{n+1}}{1-1/3}
\]
令 \(n\to\infty\)，由于 \(\left(\frac13\right)^{n+1}\to0\)，极限为 \(\frac{1}{2/3}=\frac32\). 所以填 \(\frac32\).
\end{solution}

\begin{problem}
由 \(1\)，\(2\)，\(3\)，\(4\) 四个数字组成的没有重复数字的四位数中，比 \(1234\) 大的数共有\fillinblank{}个（用具体数字作答）.
\end{problem}

\begin{answer}
\(23\)
\end{answer}

\begin{solution}
四个数字无重复组成四位数共有 \(4!=24\) 个. 首位为 \(2\)、\(3\) 或 \(4\) 的数共有 \(3\cdot3!=18\) 个，全部大于 \(1234\). 首位为 \(1\) 时，若第二位为 \(3\) 或 \(4\)，共有 \(2\cdot2!=4\) 个；第二位为 \(2\) 时只有 \(1243\) 大于 \(1234\)，而 \(1234\) 本身不计入. 所求总数为 \(18+4+1=23\)，所以填 \(23\).
\end{solution}

\begin{problem}
如果虚数 \(z\) 满足 \(z^3=8\)，那么 \(z^3+z^2+2z+2\) 的值是\fillinblank{}.
\end{problem}

\begin{answer}
\(6\)
\end{answer}

\begin{solution}
由 \(z^3=8\)，得
\(z^3-8=(z-2)(z^2+2z+4)=0\). 题设中的“虚数”不是实数，故 \(z\neq2\)，从而 \(z^2+2z+4=0\). 因此
\[
z^3+z^2+2z+2
=8+(z^2+2z+4)-2
=6
\]
所以填 \(6\).
\end{solution}

\section{解答题}

\begin{problem}
在平面 \(\beta\) 内有 \(\triangle ABC\)，在平面 \(\beta\) 外有点 \(S\)，斜线 \(SA\perp AC\)，\(SB\perp BC\). 且斜线 \(SA\)，\(SB\) 分别与平面 \(\beta\) 所成的角相等.

\begin{enumerate}
\item 求证：\(AC=BC\)；
\item 又设点 \(S\) 与平面 \(\beta\) 的距离为 \(4\text{ 厘米}\)，\(AC\perp BC\)，且 \(AB=6\text{ 厘米}\)，求点 \(S\) 与直线 \(AB\) 的距离.
\end{enumerate}
\end{problem}

\begin{solution}
\begin{enumerate}
\item 过 \(S\) 作 \(SD\perp \beta\)，\(D\) 是垂足，连结 \(AD\)，\(BD\). 则 \(SD\perp AD\)，\(SD\perp BD\)，且 \(\angle SAD\) 和 \(\angle SBD\) 分别是 \(SA\)，\(SB\) 和平面 \(\beta\) 所成的角. 依设，\(\angle SAD=\angle SBD\). 又 \(SD\) 是 \(\mathrm{Rt}\triangle SAD\) 和 \(\mathrm{Rt}\triangle SBD\) 的公共边，因此 \(\triangle SAD\cong \triangle SBD\)，所以 \(SA=SB\).
依设，\(SA\perp AC\)，\(SB\perp BC\). 连结 \(SC\). 在 \(\mathrm{Rt}\triangle SAC\) 和 \(\mathrm{Rt}\triangle SBC\) 中，\(SC=SC\)，\(SA=SB\)，故 \(\triangle SAC\cong \triangle SBC\)，所以 \(AC=BC\).

\item 由（1）的证明知 \(SA=SB\)，即 \(\triangle SAB\) 是等腰三角形. 取 \(AB\) 的中点 \(E\)，连结 \(SE\)，则 \(SE\perp AB\). 所以，\(SE\) 的长就是点 \(S\) 与直线 \(AB\) 的距离.
因为 \(SD\perp \beta\)，\(D\in \beta\)，\(SA\perp AC\)，\(AC\subset \beta\)，所以 \(AC\perp AD\)（三垂线定理的逆定理）. 同理，\(BC\perp BD\). 又 \(AC\perp BC\) 且 \(AC=BC\)，因此四边形 \(ACBD\) 是正方形.
连结 \(DE\). 因为 \(E\) 是对角线 \(AB\) 的中点，所以 \(DE=\frac{1}{2}AB\). 又因为 \(SD\perp \beta\)，\(DE\subset \beta\)，所以 \(SD\perp DE\). 依设 \(SD=4\)，且 \(AB=6\)，因此 \(DE=3\). 从而
\[
SE=\sqrt{SD^2+DE^2}=\sqrt{16+9}=5
\]
即点 \(S\) 与直线 \(AB\) 的距离为 \(5\text{ 厘米}\).
\end{enumerate}
\end{solution}

\begin{problem}
设 \(\i\) 是虚数单位，复数 \(z\) 和 \(\omega\) 满足 \(z\omega+2\i z-2\i\omega+1=0\).

\begin{enumerate}
\item 如果 \(z\) 和 \(\omega\) 又满足 \(\overline{\omega}-z=2\i\)，求 \(z\) 和 \(\omega\) 的值；
\item 求证：如果 \(\left| z \right|=\sqrt{3}\)，那么 \(\left| \omega-4\i \right|\) 的值是一个常数，并求出这个常数.
\end{enumerate}
\end{problem}

\begin{solution}
\begin{enumerate}
\item 因为 \(\overline{\omega}-z=2\i\)，所以 \(z=\overline{\omega}-2\i\). 又 \(z\omega+2\i z-2\i\omega+1=0\)，所以
\(\left( \overline{\omega}-2\i \right)\left( \omega+2\i \right)-2\i\omega+1=0\).
设 \(\omega=x+y\i\)，其中 \(x\)，\(y\) 是实数，则
\(\left[ x^2+\left( y+2 \right)^2+2y+1 \right]-2x\i=0\)，
即
\(\left( x^2+y^2+6y+5 \right)-2x\i=0\).
于是
\[
\begin{cases}
x=0,\\
y^2+6y+5=0.
\end{cases}
\]
解得 \(x=0\)，\(y=-1\)；或 \(x=0\)，\(y=-5\). 所以 \(\omega=-\i\) 或 \(\omega=-5\i\). 由于 \(z=\overline{\omega}-2\i\)，因此 \(z=-\i\) 或 \(z=3\i\). 所以所求的值为 \(z=-\i\)，\(\omega=-\i\)；或 \(z=3\i\)，\(\omega=-5\i\).

\item 因为 \(z\omega+2\i z-2\i\omega+1=0\)，所以
\(z\left( \omega+2\i \right)=2\i\omega-1\).
从而
\(\left| z \right|^2\left| \omega+2\i \right|^2=\left| 2\i\omega-1 \right|^2\).
又因为 \(\left| z \right|=\sqrt{3}\)，且
\(\left| \omega+2\i \right|^2=\left( \omega+2\i \right)\left( \overline{\omega}-2\i \right)=\omega\overline{\omega}+2\i\overline{\omega}-2\i\omega+4\)，
\(\left| 2\i\omega-1 \right|^2=\left( 2\i\omega-1 \right)\left( -2\i\overline{\omega}-1 \right)=4\omega\overline{\omega}+2\i\overline{\omega}-2\i\omega+1\).
将这三式代入前式，整理得
\(\omega\overline{\omega}-4\i\overline{\omega}+4\i\omega=11\).
因此
\[
\left| \omega-4\i \right|^2=\left( \omega-4\i \right)\left( \overline{\omega}+4\i \right)=\omega\overline{\omega}-4\i\overline{\omega}+4\i\omega+16=27
\]
所以
\(\left| \omega-4\i \right|=\sqrt{27}=3\sqrt{3}\).
这就证明了 \(\left| \omega-4\i \right|\) 的值是常数 \(3\sqrt{3}\).
\end{enumerate}
\end{solution}

\begin{problem}
已知实数 \(a>0\)，函数 \(f(x)=x^2\log_4 a\).

\begin{enumerate}
\item 如果方程 \(f(x-1)+2x=0\) 无实根，求 \(a\) 的取值范围；
\item 如果 \(a=4\)，在抛物线 \(y=f\left( \frac{x}{2} \right)\) 上有两点 \(A(x_1,y_1)\) 和 \(B(x_2,y_2)\) 满足 \(\left| AB \right|=y_1+y_2+2\)，求证：点 \(A\)，\(B\) 和这抛物线的焦点三点共线.
\end{enumerate}
\end{problem}

\begin{solution}
\begin{enumerate}
\item 因为 \(f(x)=x^2\log_4 a\)，所以
\(f(x-1)=\left( x-1 \right)^2\log_4 a=\left( x^2-2x+1 \right)\log_4 a\).
原方程即为
\(x^2\log_4 a+2\left( 1-\log_4 a \right)x+\log_4 a=0\).
要使方程无实根，必须且只须判别式
\(4\left( 1-\log_4 a \right)^2-4\log_4^2 a<0\).
即
\(1-2\log_4 a<0\)，
所以
\(\log_4 a>\frac{1}{2}\).
因为对数的底 \(4>1\)，由对数函数的单调性，得 \(a>4^{\frac{1}{2}}=2\). 将 \(a>2\) 代入检验，确是原不等式的解，因此 \(a>2\) 为所求的 \(a\) 的取值范围.

\item 因为 \(a=4\)，所以
\(f\left( \frac{x}{2} \right)=\left( \frac{x}{2} \right)^2\log_4 4=\frac{x^2}{4}\).
于是抛物线 \(y=f\left( \frac{x}{2} \right)\) 即为
\(x^2=4y\)，
它的焦点为 \(F(0,1)\)，准线 \(l\) 的方程为 \(y=-1\). 点 \(A(x_1,y_1)\) 和 \(B(x_2,y_2)\) 在抛物线上，故 \(y_1=\frac{1}{4}x_1^2\geqslant 0\)，\(y_2=\frac{1}{4}x_2^2\geqslant 0\). 且 \(A\)，\(B\) 分别与准线 \(l\) 的距离为 \(d_1=y_1+1\) 和 \(d_2=y_2+1\).

\solutionfigure{\bitmapfigure[max width=0.55\linewidth]{img/1989/guangdong_science/q4_solution_fig1.png}}

根据抛物线的几何性质，\(\left| AF \right|=d_1\)，\(\left| BF \right|=d_2\). 所以
\(\left| AF \right|+\left| BF \right|=y_1+y_2+2\).
依设，\(\left| AB \right|=y_1+y_2+2\)，因此
\(\left| AF \right|+\left| BF \right|=\left| AB \right|\).
如果 \(A\)，\(B\)，\(F\) 三点不共线，那么考察 \(\triangle ABF\)，根据三角形的性质，有 \(\left| AF \right|+\left| BF \right|>\left| AB \right|\)，这与上式矛盾. 所以，\(A\)，\(B\) 和 \(F\) 三点必共线.
\end{enumerate}
\end{solution}

\begin{problem}
设圆 \(C\) 的方程为 \(x^2+y^2-2x\left( \frac{1-\cos \theta}{1+\cos \theta} \right)-2y\tan \frac{\theta}{2}+\left( \frac{1-\cos \theta}{1+\cos \theta} \right)^2=0\)，式中 \(\theta\) 是实数，且 \(0<\theta<\pi\).

\begin{enumerate}
\item 当 \(\theta\) 在区间 \((0,\pi)\) 内变动时，求圆 \(C\) 的圆心轨迹的方程（写成普通形式），并画出轨迹曲线的草图；
\item 又设 \(\theta_1\)，\(\theta_2\)，\(\theta_3\) 都是区间 \((0,\pi)\) 内的实数，且 \(\theta_1\)，\(\theta_2\)，\(\theta_3\) 成为公差不为零的等差数列，当 \(\theta\) 依次取值 \(\theta_1\)，\(\theta_2\) 或 \(\theta_3\) 时，所对应的圆 \(C\) 的半径依次为 \(r_1\)，\(r_2\) 和 \(r_3\)，试问：\(r_1\)，\(r_2\)，\(r_3\) 是否成等比数列？为什么？
\end{enumerate}
\end{problem}

\begin{solution}
\begin{enumerate}
\item 化圆 \(C\) 的方程为 \(\left( x-\frac{1-\cos \theta}{1+\cos \theta} \right)^2+\left( y-\tan \frac{\theta}{2} \right)^2=\tan^2 \frac{\theta}{2}\).
得圆 \(C\) 的圆心坐标为
\[
\begin{cases}
x=\frac{1-\cos \theta}{1+\cos \theta},\\
y=\tan \frac{\theta}{2},
\end{cases}
\qquad 0<\theta<\pi
\]
又
\(\frac{1-\cos \theta}{1+\cos \theta}=\frac{2\sin^2 \frac{\theta}{2}}{2\cos^2 \frac{\theta}{2}}=\tan^2 \frac{\theta}{2}\)，
因此所求的圆心轨迹的普通方程为
\(y^2=x\)（\(y>0\)）.
也可写成 \(y=\sqrt{x}\)（\(x>0\)）.

\solutionfigure{\bitmapfigure[max width=0.50\linewidth]{img/1989/guangdong_science/q5_solution_fig1.png}}

这表明圆心轨迹曲线是位于第一象限中的半支抛物线，不含坐标原点.

\item
\begin{enumerate}
\item 由圆 \(C\) 的方程可知圆 \(C\) 的半径为
\(r=\left| \tan \frac{\theta}{2} \right|\).
因为 \(0<\theta<\pi\)，所以 \(r=\tan \frac{\theta}{2}\). 依设
\(r_1=\tan \frac{\theta_1}{2}\)，\(r_2=\tan \frac{\theta_2}{2}\)，\(r_3=\tan \frac{\theta_3}{2}\).
又因为 \(\theta_1\)，\(\theta_2\)，\(\theta_3\) 是公差不为零的等差数列，所以
\(\theta_3-\theta_2=\theta_2-\theta_1\neq 0\)，
即
\(\frac{\theta_1}{2}+\frac{\theta_3}{2}=\theta_2\).

先看 \(\theta_2=\frac{\pi}{2}\) 的情形. 这时
\(r_2=\tan \frac{\pi}{4}=1\).
又由 \(\frac{\theta_1}{2}+\frac{\theta_3}{2}=\frac{\pi}{2}\)，且 \(0<\frac{\theta_1}{2}<\frac{\pi}{2}\)，\(0<\frac{\theta_2}{2}<\frac{\pi}{2}\)，得
\(\tan \frac{\theta_1}{2}\tan \frac{\theta_3}{2}=1\).
因此
\(r_1r_3=1=r_2^2\).
并且 \(r_1\)，\(r_2\)，\(r_3\) 都不为零，所以这时 \(r_1\)，\(r_2\)，\(r_3\) 成等比数列.

\item 再看 \(\theta_2\neq \frac{\pi}{2}\) 的情形. 由 \(\frac{\theta_1}{2}+\frac{\theta_3}{2}=\theta_2\)，可得
\(\frac{\tan \frac{\theta_1}{2}+\tan \frac{\theta_3}{2}}{1-\tan \frac{\theta_1}{2}\tan \frac{\theta_3}{2}}=\frac{2\tan \frac{\theta_2}{2}}{1-\tan^2 \frac{\theta_2}{2}}\).
于是
\(\frac{r_1+r_3}{1-r_1r_3}=\frac{2r_2}{1-r_2^2}\).
如果 \(r_1\)，\(r_2\)，\(r_3\) 成等比数列，则 \(r_1r_3=r_2^2\)，故 \(1-r_1r_3=1-r_2^2\). 因此前式化为
\(r_1+r_3=2r_2\).
于是
\(r_1^2+2r_1r_3+r_3^2=4r_2^2\).
再由 \(r_1r_3=r_2^2\)，得
\(r_1^2-2r_1r_3+r_3^2=0\)，
即
\(\left( r_1-r_3 \right)^2=0\).
所以 \(r_1=r_3\). 又因为 \(r_1+r_3=2r_2\)，所以 \(r_1=r_2=r_3\). 这就是说
\(\tan \frac{\theta_1}{2}=\tan \frac{\theta_2}{2}=\tan \frac{\theta_3}{2}\).
由于 \(\frac{\theta_1}{2}\)，\(\frac{\theta_2}{2}\)，\(\frac{\theta_3}{2}\) 都在区间 \(\left( 0,\frac{\pi}{2} \right)\) 内，上式成立必有
\(\frac{\theta_1}{2}=\frac{\theta_2}{2}=\frac{\theta_3}{2}\).
于是 \(\theta_1=\theta_2=\theta_3\)，这与它们的公差不为零矛盾. 所以当 \(\theta_2\neq \frac{\pi}{2}\) 时，\(r_1\)，\(r_2\)，\(r_3\) 不可能成为等比数列.

\end{enumerate}
综上所述，当 \(\theta_2=\frac{\pi}{2}\) 时，\(r_1\)，\(r_2\)，\(r_3\) 必定成为等比数列；当 \(\theta_2\neq \frac{\pi}{2}\) 时，\(r_1\)，\(r_2\)，\(r_3\) 不可能成为等比数列.
\end{enumerate}
\end{solution}
